% GENERATED FILE. Edit canonical lessons, then run npm run generate:books.
% canonical-source-hash: fnv1a64:ac37732d057338be
% canonical-lessons: KA-C134-ili, KA-C134-emme, KA-C134-kuri, KA-C134-mukha, KA-C134-nalige

\chapter{Rat, Buffalo, Sheep, Face, Tongue}
\label{ch:ka-rat-buffalo-sheep-face-tongue}
\hlchaptermodality{134}

\begin{chapteropening}
\textbf{By the end of this chapter:} \emph{I can say a rat, a buffalo, a sheep, the face and the tongue.}

Complete the last lesson of chapter 134: I can say a rat, a buffalo, a sheep, the face and the tongue.
\end{chapteropening}

\section[\texorpdfstring{\kn{ಇಲಿ} (ili)}{ಇಲಿ (ili)}]{\kn{ಇಲಿ} (ili) — a rat, a mouse}

\label{lesson:KA-C134-ili}



\begin{quote}
Before the new one: say the Kannada for a donkey, then the Kannada for a camel.
\end{quote}

\subsection*{You'll want to know: \kn{ಇಲಿ}}
\textbf{\kn{ಇಲಿ}} — \emph{ili} — \textquotedblleft{}a rat, a mouse\textquotedblright{}.

\begin{grammarlens}[title={What you've built}]
One more word for this chapter.
\end{grammarlens}

\subsection*{Your turn}
\begin{itemize}
\raggedright
  \item \emph{Say it:} \emph{ili}
  \item \emph{Say it:} \emph{ili}, once more
  \item \emph{Say it:} say \emph{ili}
  \item \emph{Recall:} say the Kannada for a donkey, then the Kannada for a camel, then say \emph{ili} again
\end{itemize}

\begin{tcolorbox}[breakable,colback=teal!4,colframe=teal!35!black,title={Before you move on}]
What does \kn{ಇಲಿ} mean? (\textquotedblleft{}A rat, a mouse\textquotedblright{}.) Say it once more.
\end{tcolorbox}
\section[\texorpdfstring{\kn{ಎಮ್ಮೆ} (emme)}{ಎಮ್ಮೆ (emme)}]{\kn{ಎಮ್ಮೆ} (emme) — a buffalo}

\label{lesson:KA-C134-emme}



\begin{quote}
Before the new one: say the Kannada for a camel, then the Kannada for a rat.
\end{quote}

\subsection*{You'll want to know: \kn{ಎಮ್ಮೆ}}
\textbf{\kn{ಎಮ್ಮೆ}} — \emph{emme} — \textquotedblleft{}a buffalo\textquotedblright{}.

\begin{grammarlens}[title={What you've built}]
One more word for this chapter.
\end{grammarlens}

\subsection*{Your turn}
\begin{itemize}
\raggedright
  \item \emph{Say it:} \emph{emme}
  \item \emph{Say it:} \emph{emme}, once more
  \item \emph{Say it:} say \emph{emme}
  \item \emph{Recall:} say the Kannada for a camel, then the Kannada for a rat, then say \emph{emme} again
\end{itemize}

\begin{tcolorbox}[breakable,colback=teal!4,colframe=teal!35!black,title={Before you move on}]
What does \kn{ಎಮ್ಮೆ} mean? (\textquotedblleft{}A buffalo\textquotedblright{}.) Say it once more.
\end{tcolorbox}
\section[\texorpdfstring{\kn{ಕುರಿ} (kuri)}{ಕುರಿ (kuri)}]{\kn{ಕುರಿ} (kuri) — a sheep}

\label{lesson:KA-C134-kuri}



\begin{quote}
Before the new one: say the Kannada for a rat, then the Kannada for a buffalo.
\end{quote}

\subsection*{You'll want to know: \kn{ಕುರಿ}}
\textbf{\kn{ಕುರಿ}} — \emph{kuri} — \textquotedblleft{}a sheep\textquotedblright{}.

\begin{grammarlens}[title={What you've built}]
One more word for this chapter.
\end{grammarlens}

\subsection*{Your turn}
\begin{itemize}
\raggedright
  \item \emph{Say it:} \emph{kuri}
  \item \emph{Say it:} \emph{kuri}, once more
  \item \emph{Say it:} say \emph{kuri}
  \item \emph{Recall:} say the Kannada for a rat, then the Kannada for a buffalo, then say \emph{kuri} again
\end{itemize}

\begin{tcolorbox}[breakable,colback=teal!4,colframe=teal!35!black,title={Before you move on}]
What does \kn{ಕುರಿ} mean? (\textquotedblleft{}A sheep\textquotedblright{}.) Say it once more.
\end{tcolorbox}
\section[\texorpdfstring{\kn{ಮುಖ} (mukha)}{ಮುಖ (mukha)}]{\kn{ಮುಖ} (mukha) — the face}

\label{lesson:KA-C134-mukha}



\begin{quote}
Before the new one: say the Kannada for a buffalo, then the Kannada for a sheep.
\end{quote}

\subsection*{You'll want to know: \kn{ಮುಖ}}
\textbf{\kn{ಮುಖ}} — \emph{mukha} — \textquotedblleft{}the face\textquotedblright{}.

\begin{grammarlens}[title={What you've built}]
One more word for this chapter.
\end{grammarlens}

\subsection*{Your turn}
\begin{itemize}
\raggedright
  \item \emph{Say it:} \emph{mukha}
  \item \emph{Say it:} \emph{mukha}, once more
  \item \emph{Say it:} say \emph{mukha}
  \item \emph{Recall:} say the Kannada for a buffalo, then the Kannada for a sheep, then say \emph{mukha} again
\end{itemize}

\begin{tcolorbox}[breakable,colback=teal!4,colframe=teal!35!black,title={Before you move on}]
What does \kn{ಮುಖ} mean? (\textquotedblleft{}The face\textquotedblright{}.) Say it once more.
\end{tcolorbox}
\section[\texorpdfstring{\kn{ನಾಲಿಗೆ} (nālige)}{ನಾಲಿಗೆ (nālige)}]{\kn{ನಾಲಿಗೆ} (nālige) — the tongue}

\label{lesson:KA-C134-nalige}



\begin{quote}
Before the new one: say the Kannada for a sheep, then the Kannada for the face.
\end{quote}

\subsection*{You'll want to know: \kn{ನಾಲಿಗೆ}}
\textbf{\kn{ನಾಲಿಗೆ}} — \emph{nālige} — \textquotedblleft{}the tongue\textquotedblright{}.

\begin{grammarlens}[title={What you've built}]
One more word for this chapter.
\end{grammarlens}

\subsection*{Your turn}
\begin{itemize}
\raggedright
  \item \emph{Say it:} \emph{nālige}
  \item \emph{Say it:} \emph{nālige}, once more
  \item \emph{Say it:} say \emph{nālige}
  \item \emph{Recall:} say the Kannada for a sheep, then the Kannada for the face, then say \emph{nālige} again
\end{itemize}

\begin{tcolorbox}[breakable,colback=teal!4,colframe=teal!35!black,title={Before you move on}]
What does \kn{ನಾಲಿಗೆ} mean? (\textquotedblleft{}The tongue\textquotedblright{}.) Say it once more.
\end{tcolorbox}
